\documentclass[../Main.tex]{subfiles}
\begin{document}

Chương 4 đã trình bày tổng thể kiến trúc phần mềm, thiết kế chi tiết các gói, cơ sở dữ liệu cùng quy trình triển khai của hệ thống VoltGO. Trên cơ sở đó, Chương 5 tập trung trình bày những đóng góp kỹ thuật có giá trị nhất mà đồ án đã đạt được. Mỗi đóng góp được tổ chức thành một mục độc lập, gồm ba phần: phát biểu bài toán, giải pháp đề xuất và kết quả thực nghiệm. Phần 5.1 trình bày tính năng dẫn đường nhận biết xe điện kết hợp với gợi ý trạm sạc dọc tuyến, phần 5.2 trình bày cơ chế xếp hạng đa yếu tố cho trạm ứng viên, phần 5.3 mô tả bộ lọc an toàn về mức pin tới trạm, phần 5.4 trình bày quy trình gợi ý thời điểm sạc thông minh kết hợp lịch sử đặt chỗ, và phần 5.5 tổng hợp lại những điểm đã đạt được của nhóm tính năng này.

\section{Tính năng dẫn đường nhận biết xe điện kết hợp gợi ý trạm sạc dọc tuyến}
\label{section:5.1}

\subsection{Phát biểu bài toán}

Use case UC-AI-003 đã được trình bày trong mục \ref{subsection:2.3.3} cho thấy người lái xe điện cần nhận được danh sách trạm sạc được gợi ý khi lập kế hoạch cho một hành trình. Một giải pháp ngây thơ là tìm tất cả các trạm đã công bố trong bán kính nào đó tính từ vị trí hiện tại, nhưng cách làm này có hai hạn chế nghiêm trọng. Thứ nhất, các trạm được trả về có thể nằm hoàn toàn ngoài hướng di chuyển thực tế của người dùng, buộc họ phải rẽ ngược hoặc đi vòng xa hơn quãng đường còn lại. Thứ hai, cách tiếp cận theo bán kính không tận dụng được thông tin tuyến đường mà máy chủ đã trả về từ dịch vụ định tuyến, dẫn đến việc người dùng nhận được các gợi ý không khả thi về mặt hành trình. Bài toán đặt ra đòi hỏi hệ thống phải kết hợp được hai nguồn dữ liệu không đồng nhất: một tuyến đường đa điểm do dịch vụ định tuyến bên ngoài cung cấp, và tập các trạm sạc đã công bố được lưu trữ trong cơ sở dữ liệu có hỗ trợ truy vấn không gian, sao cho danh sách kết quả phải vừa nằm dọc theo tuyến đường thực tế, vừa phù hợp với tình trạng pin và phạm vi hoạt động còn lại của xe.

\subsection{Giải pháp}

Hệ thống VoltGO giải quyết bài toán trên bằng một pipeline gồm bốn bước liên tiếp, được cài đặt trong lớp dịch vụ \texttt{RoutingService} ở phía máy chủ. Đầu tiên, khi người dùng chọn một điểm đến, dịch vụ gọi máy chủ OSRM công cộng tại \texttt{router.project-osrm.org} thông qua \texttt{WebClient} của Spring để lấy về cung đường tối ưu dưới dạng chuỗi các cặp tọa độ (kinh độ, vĩ độ) cùng với tổng quãng đường và thời gian di chuyển ước tính. Trước mỗi lần gọi, hệ thống kiểm tra tính hợp lệ của tọa độ đầu vào dựa trên khoảng giá trị [-90, 90] cho vĩ độ và [-180, 180] cho kinh độ, đồng thời yêu cầu OSRM trả về cung đường đầy đủ (tham số \texttt{overview=full}) với định dạng hình học GeoJSON, nhằm phục vụ cho cả việc hiển thị trên bản đồ lẫn tính toán hành lang tìm kiếm. Trong trường hợp lần gọi đầu tiên thất bại vì lỗi tạm thời của OSRM hoặc sự cố mạng, dịch vụ tự động thử lại tối đa một lần theo cơ chế retry với mã lỗi \texttt{INTERNAL\_ERROR} hoặc \texttt{SERVICE\_UNAVAILABLE}, đảm bảo tính sẵn sàng trước những gián đoạn không thường xuyên.

Sau khi nhận được cung đường từ OSRM, hệ thống chuyển đổi chuỗi đa điểm này sang định dạng WKT (Well-Known Text) \texttt{LINESTRING} để có thể sử dụng các hàm không gian của PostGIS. Bước thứ hai là tìm kiếm các trạm nằm trong một hành lang quanh tuyến đường. Khoảng cách mặc định của hành lang được cấu hình trong \texttt{application.yml} với giá trị 5000 mét, tức là hệ thống sẽ xem xét tất cả các trạm cách cung đường không quá năm kilômét. Câu truy vấn PostGIS sử dụng \texttt{ST\_DWithin} với kiểu \texttt{geography} để tính khoảng cách chính xác trên bề mặt cầu trái đất, đồng thời kết hợp \texttt{ST\_Distance} để lấy khoảng cách từ tuyến đến mỗi trạm và khoảng cách từ điểm xuất phát đến trạm. Nhờ chỉ mục không gian GIST trên cột \texttt{location} của bảng \texttt{station\_version}, truy vấn này có độ phức tạp gần $O(\log n)$ thay vì phải quét toàn bộ bảng, đảm bảo thời gian phản hồi nhanh ngay cả khi số lượng trạm tăng cao.

Một điểm đặc biệt của giải pháp là cơ chế \emph{fallback theo bậc} cho bán kính hành lang. Khi không tìm được trạm nào trong bán kính 5000 mét, hệ thống tự động thử lại với các bán kính lớn hơn được liệt kê trong cấu hình \texttt{fallback-buffer-meters}, mặc định là 1000, 3000, 5000, rồi 10000 mét. Cách làm này giải quyết tình huống người dùng đi vào vùng có mật độ trạm thưa, đảm bảo rằng hệ thống vẫn có khả năng đưa ra gợi ý thay vì trả về danh sách rỗng. Toàn bộ quá trình này được ghi nhận vào nhật ký với một mã truy vết ngắn tám ký tự sinh ngẫu nhiên cho mỗi yêu cầu, giúp việc truy vết sự cố trong môi trường nhiều người dùng trở nên khả thi. Bước thứ ba của pipeline là tính toán phạm vi hoạt động còn lại dựa trên thông tin pin do ứng dụng khách gửi lên. Hệ thống ước tính quãng đường còn đi được bằng tích của phần trăm pin hiện tại và phạm vi tối đa ở mức pin đầy, sau đó so sánh với tổng quãng đường của tuyến cộng thêm một ngưỡng an toàn cấu hình được (mặc định 5 phần trăm phạm vi xe) để xác định người dùng có cần dừng sạc hay không.

Bước cuối cùng là tổng hợp kết quả. Hệ thống chọn ra tối đa ba trạm tốt nhất theo thứ tự điểm số tăng dần do cơ chế chấm điểm ở phần 5.2 tính ra, đánh dấu trạm đứng đầu là điểm dừng tối ưu và đánh dấu cả ba là các trạm được khuyến nghị, đồng thời gửi kèm theo cung đường, hộp bao quanh, và chuỗi lý do bằng ngôn ngữ tự nhiên giải thích vì sao trạm đó được xếp hạng như vậy. Phía máy khách chỉ đơn thuần hiển thị các trường này lên bản đồ, vẽ polyline cung đường và các biểu tượng trạm với màu sắc phân biệt giữa trạm tối ưu (xanh lá) và trạm được khuyến nghị (cam), hoàn toàn không tham gia vào tính toán. Cách phân chia trách nhiệm rõ ràng theo mô hình máy chủ làm hết, máy khách chỉ hiển thị giúp cho thuật toán có thể được điều chỉnh và tối ưu mà không cần phát hành lại ứng dụng khách.

\begin{figure}[H]
    \centering
    \includegraphics[width=0.9\linewidth]{Figure/Figure_Placeholder.png}
    \caption{Luồng xử lý dẫn đường nhận biết xe điện kết hợp gợi ý trạm sạc dọc tuyến}
    \label{fig:routing_pipeline}
\end{figure}

\subsection{Kết quả đạt được}

Hệ thống đã được tích hợp đầy đủ vào ứng dụng di động dành cho người lái xe điện và hoạt động ổn định trong các tình huống thực tế. Người dùng nhập điểm đến thông qua thanh tìm kiếm với gợi ý địa điểm hoặc nhấn giữ trên bản đồ, sau đó cung cấp mức pin hiện tại cùng phạm vi tối đa của xe. Kết quả trả về bao gồm cung đường tối ưu vẽ trực tiếp lên bản đồ, biểu tượng vị trí điểm đến, các trạm sạc nằm dọc tuyến được đánh dấu rõ ràng, cùng một bảng thông tin chi tiết về quãng đường, thời gian di chuyển, trạng thái pin còn lại và trạm dừng tối ưu. Trong trường hợp quãng đường còn lại vượt quá phạm vi pin dự kiến, hệ thống sẽ tự động kích hoạt cờ "cần sạc" và ưu tiên hiển thị các trạm gợi ý; ngược lại nếu pin đủ, hệ thống hiển thị thông báo "đủ cho chuyến đi" để người dùng yên tâm. Khi xảy ra lỗi mạng hoặc dịch vụ OSRM không phản hồi sau lần thử lại, giao diện sẽ hiển thị một khung lỗi với nút nhấn cho phép thử lại mà không cần phải nhập lại điểm đến, giúp giảm thiểu sự bất tiện cho người dùng. Việc tách biệt phần tính toán hoàn toàn về phía máy chủ cũng giúp hệ thống có thể dễ dàng tinh chỉnh trọng số, thay đổi chiến lược fallback hoặc bổ sung thêm các yếu tố mới vào hàm chấm điểm mà không phải phát hành lại ứng dụng trên các kho ứng dụng.

\section{Cơ chế xếp hạng đa yếu tố cho trạm sạc ứng viên}
\label{section:5.2}

\subsection{Phát biểu bài toán}

Sau khi bước tìm kiếm không gian ở phần 5.1 thu hẹp tập ứng viên xuống còn một số lượng vừa phải, bài toán tiếp theo là xếp hạng các trạm này sao cho trạm đứng đầu phải là lựa chọn tốt nhất theo quan điểm của người dùng. Hệ thống cần cân nhắc đồng thời nhiều yếu tố có đơn vị đo lường khác nhau: khoảng cách phải đi thêm so với tuyến gốc được tính bằng mét, công suất sạc tối đa tính bằng kilôwatt, số cổng đang rảnh là số nguyên, thời gian chờ và thời gian sạc ước tính tính bằng phút, và điểm tin cậy là số thực trong khoảng từ 0 đến 100. Việc kết hợp trực tiếp các đại lượng này bằng phép cộng đơn thuần là không thể vì chúng có thang đo khác nhau; nếu cố gắng đưa về cùng đơn vị thì sẽ mất đi ý nghĩa vật lý của từng yếu tố. Hơn nữa, đặc thù của trạm đổi pin còn khác hẳn so với trạm sạc truyền thống: trạm đổi pin có thời gian phục vụ cố định khoảng năm phút và phụ thuộc vào số pin có sẵn chứ không phải công suất sạc, do đó cần một hàm chấm điểm riêng để phản ánh đúng bản chất của từng loại hình dịch vụ.

\subsection{Giải pháp}

Giải pháp của đồ án là xây dựng một hàm chấm điểm thống nhất mà trong đó điểm càng thấp thì trạm càng tốt, cho phép so sánh trực tiếp giữa các trạm thuộc cùng loại hình dịch vụ lẫn khác loại. Hàm chấm điểm này được cài đặt trong phương thức \texttt{computeUnifiedScore} của lớp \texttt{StationQueryRepositoryImpl} và có hai biến thể: một biến thể đầy đủ khi người dùng cung cấp thông tin pin và phạm vi xe, và một biến thể đơn giản \texttt{computeUnifiedScoreNoEv} cho trường hợp ngược lại.

Đối với trạm sạc thuần túy, điểm số được tính theo công thức tổng của năm thành phần, mỗi thành phần là tích của một trọng số không đổi và một đại lượng đầu vào. Thành phần thứ nhất phạt khoảng cách đi vòng với trọng số 2,0, theo nguyên tắc mỗi mét đi vòng đều tốn thời gian và năng lượng. Thành phần thứ hai thưởng cho công suất sạc cao với trọng số -1,5, nghĩa là công suất lớn hơn sẽ làm giảm điểm và cải thiện thứ hạng. Thành phần thứ ba thưởng cho số cổng đang rảnh với trọng số -5,0, vì nhiều cổng rảnh đồng nghĩa với xác suất phải chờ thấp hơn. Thành phần thứ bốn phạt thời gian chờ ước tính, được tính bằng 10 phút khi không còn cổng nào rảnh và bằng 0 trong trường hợp ngược lại, với trọng số 1,5. Thành phần thứ năm phạt thời gian sạc ước tính, được tính dựa trên phần năng lượng cần bổ sung chia cho công suất hiệu dụng, với trọng số 0,3. Trường hợp trạm không có công suất sạc khai báo sẽ bị phạt bằng hằng số \texttt{PENALTY\_UNREACHABLE} có giá trị 100 nghìn, lớn hơn gấp nhiều lần bất kỳ điểm hợp lệ nào, đảm bảo trạm đó luôn bị xếp cuối.

Đối với trạm đổi pin, hàm chấm điểm hoàn toàn khác, phản ánh đặc thù dịch vụ. Nếu trạm không còn pin nào sẵn sàng, trạm đó bị gán điểm \texttt{SWAP\_PENALTY\_NO\_BATTERY} bằng 80 nghìn, đảm bảo bị loại khỏi danh sách khuyến nghị. Nếu còn pin, điểm được tính bằng tổng của ba thành phần: khoảng cách đi vòng với trọng số 2,0, số pin sẵn có với trọng số -3,0 (thưởng cho trạm có nhiều pin hơn) và giá cơ bản với trọng số 0,01. Sở dĩ hệ số của giá nhỏ đến vậy là vì đơn vị của giá là đồng Việt Nam, giá trị lớn hơn nhiều so với khoảng cách tính bằng mét, nên cần trọng số nhỏ để giá không áp đảo các yếu tố khác.

Đối với trạm hỗ trợ cả hai dịch vụ sạc và đổi pin, hệ thống tính điểm theo cả hai cách rồi lấy giá trị nhỏ hơn, phản ánh triết lý người dùng có thể chọn cách nào cũng được và hệ thống nên ưu tiên phương án cho điểm tốt hơn. Khi không có thông số pin của xe, biến thể không có tham số EV sử dụng công thức đơn giản hơn với phần phạt thời gian sạc cố định ba mươi phút thay cho giá trị tính toán theo năng lượng cần bổ sung, đảm bảo hệ thống vẫn cho kết quả hợp lý trong trường hợp người dùng chưa khai báo đặc điểm xe. Ngoài giá trị số, mỗi trạm còn được gán kèm một chuỗi giải thích bằng ngôn ngữ tự nhiên, ví dụ như "Ngay trên tuyến, công suất 120 kW DC, ba cổng rảnh, sạc khoảng 25 phút, pin còn khoảng 30 phần trăm khi tới", giúp người dùng hiểu được lý do xếp hạng mà không cần phải diễn giải ngược các con số trừu tượng.

\subsection{Kết quả đạt được}

Hàm chấm điểm đã được tích hợp vào cả hai luồng nghiệp vụ là gợi ý trạm dọc tuyến đường và gợi ý trạm trong bán kính quanh vị trí hiện tại. Khi người dùng mở ứng dụng, bản đồ sẽ hiển thị các trạm lân cận cùng điểm tin cậy và số cổng rảnh, còn khi họ nhập điểm đến và lập kế hoạch hành trình, bản đồ sẽ hiển thị đồng thời cung đường và các trạm được khuyến nghị, với trạm xếp hạng cao nhất được đánh dấu bằng biểu tượng ngôi sao màu xanh lá cây và các trạm còn lại được đánh dấu bằng biểu tượng tia chớp màu cam. Khi người dùng nhấn vào biểu tượng của một trạm, một bảng thông tin chi tiết sẽ hiện ra cho thấy đầy đủ khoảng cách đi vòng, công suất, số cổng rảnh, thời gian sạc dự kiến, điểm tin cậy và cả chuỗi giải thích lý do. Thiết kế này giúp người dùng vừa có thể so sánh nhanh các lựa chọn thông qua vị trí và màu sắc trên bản đồ, vừa có thể đi sâu vào chi tiết từng trạm khi cần ra quyết định cuối cùng. Việc đồng nhất cách chấm điểm giữa hai luồng sử dụng cũng giúp trải nghiệm người dùng nhất quán: cùng một logic đánh giá được áp dụng cho dù họ đang tìm trạm gần nhà hay trạm dọc đường đi công tác.

\section{Bộ lọc an toàn về mức pin tới trạm}
\label{section:5.3}

\subsection{Phát biểu bài toán}

Một trạm sạc có thể nằm rất gần tuyến đường và có công suất cao nhưng vẫn trở nên vô dụng nếu người dùng không thể tới được nó vì lý do năng lượng. Ví dụ, một trạm cách điểm xuất phát 50 kilômét về hướng ngược lại với đích đến, dù gần về mặt không gian, có thể khiến người dùng phải tiêu hao lượng pin vượt quá mức có thể, đặc biệt khi mức pin ban đầu không cao hoặc phạm vi tối đa của xe thấp. Vấn đề này đặc biệt nghiêm trọng với hệ thống gợi ý theo hành lang vì nhiều trạm cùng lúc có thể "gần" theo nghĩa khoảng cách tới tuyến nhỏ nhưng "xa" theo nghĩa quãng đường thực tế từ điểm xuất phát. Nếu hệ thống đề xuất một trạm mà người dùng không thể tới vì lý do năng lượng, toàn bộ lợi ích của việc gợi ý sẽ bị triệt tiêu, thậm chí còn gây hại vì người dùng có thể đi theo hướng đó rồi mới phát hiện ra không đủ pin.

\subsection{Giải pháp}

Để giải quyết bài toán này, hệ thống tích hợp một bộ lọc an toàn trong chính phương thức \texttt{findStationsAlongRoute} của lớp \texttt{StationQueryRepositoryImpl}. Với mỗi trạm ứng viên, trước khi tính điểm, hệ thống ước tính mức pin khi tới trạm dựa trên công thức tuyến tính đơn giản: khoảng cách thực tế từ điểm xuất phát đến trạm (tính bằng PostGIS \texttt{ST\_Distance} với kiểu \texttt{geography}) nhân với mức tiêu thụ năng lượng trung bình, mặc định 0,18 kWh cho mỗi kilômét, chia cho phạm vi tối đa của xe ở mức pin đầy, nhân với 100 để ra phần trăm pin bị tiêu hao. Mức pin tới trạm bằng mức pin hiện tại trừ đi phần trăm bị tiêu hao, với giới hạn dưới bằng 0 để tránh kết quả âm.

Nếu mức pin tới trạm thấp hơn ngưỡng an toàn \texttt{min-battery-percent-at-arrival}, mặc định là 5 phần trăm, hệ thống ngay lập tức gán cho trạm đó điểm \texttt{SCORE\_UNREACHABLE} bằng 50 nghìn, một giá trị lớn hơn gấp nhiều lần bất kỳ điểm hợp lệ nào nhưng lại nhỏ hơn \texttt{PENALTY\_UNREACHABLE} 100 nghìn dành cho trạm không có công suất. Sự phân biệt này cho phép giao diện người dùng có thể hiển thị trạm không tới được với trạng thái "không thể tới, không đủ pin" khác với trạng thái "không có công suất sạc", mặc dù cả hai đều bị xếp cuối cùng. Đồng thời, lý do bằng ngôn ngữ tự nhiên cũng được thay đổi thành "Không thể tới vì không đủ pin cho quãng đường đi vòng" để giúp người dùng hiểu được tại sao trạm đó bị loại khỏi danh sách khuyến nghị.

Trong trường hợp mức pin tới trạm đạt yêu cầu, hệ thống tiếp tục tính toán năng lượng cần bổ sung cho phần còn lại của hành trình, rồi ước tính thời gian sạc bằng công thức năng lượng cần bổ sung chia cho công suất hiệu dụng (lấy giá trị nhỏ hơn giữa công suất cổng sạc và công suất sạc tối đa của xe) nhân với 60 phút, với giới hạn trên 120 phút để tránh hiển thị thời gian sạc quá dài cho những trạm có công suất thấp. Thời gian sạc ước tính này được sử dụng làm đầu vào cho hàm chấm điểm đã trình bày ở phần 5.2.

\subsection{Kết quả đạt được}

Bộ lọc an toàn đã chứng minh vai trò quan trọng trong các tình huống thực tế có độ chênh lệch lớn giữa khoảng cách tới tuyến và khoảng cách từ điểm xuất phát. Khi một người dùng ở Hà Nội lập kế hoạch đi Hải Phòng với mức pin 30 phần trăm trên một chiếc xe có phạm vi tối đa 300 kilômét, hệ thống sẽ tự động loại bỏ các trạm nằm ở hướng ngược lại về phía nam hoặc quá xa về phía bắc vì không thể đảm bảo mức pin tối thiểu khi tới, đồng thời giữ lại các trạm nằm dọc hướng Hà Nội - Hải Phòng. Nhờ vậy, người dùng không bị phân tâm bởi các gợi ý hấp dẫn về mặt thị giác nhưng không khả thi về mặt năng lượng. Trong trường hợp người dùng cố tình muốn xem các trạm bị loại, hệ thống vẫn có thể trả về chúng với cờ \texttt{isRecommended} bằng false và chuỗi giải thích kèm theo, giúp người dùng có đầy đủ thông tin để đưa ra quyết định cuối cùng.

\section{Gợi ý thời điểm sạc thông minh kết hợp lịch sử đặt chỗ}
\label{section:5.4}

\subsection{Phát biểu bài toán}

Ngoài việc chọn trạm nào để sạc, người dùng còn phải đối mặt với câu hỏi sạc vào lúc nào để tránh phải chờ lâu hoặc giảm chi phí. Ở các trạm sạc công cộng, tình trạng tắc nghẽn thường xảy ra vào giờ cao điểm buổi tối khi nhiều người đi làm về, trong khi buổi sáng sớm hoặc giữa trưa thường vắng hơn. Nếu hệ thống chỉ đơn thuần đề xuất thời điểm sạc ngẫu nhiên hoặc theo giờ mặc định thì sẽ bỏ lỡ cơ hội giúp người dùng tiết kiệm thời gian chờ. Mặt khác, thói quen sạc của mỗi người dùng cũng khác nhau: người thích sạc buổi tối, người lại thích sạc buổi sáng trước khi đi làm. Hệ thống cần cân bằng giữa việc đề xuất khung giờ ít tải cho trạm và tôn trọng thói quen của người dùng.

\subsection{Giải pháp}

Hệ thống VoltGO cung cấp một endpoint gợi ý thời điểm sạc thông minh, được cài đặt trong lớp \texttt{EvUserAiService} với phương thức \texttt{getSmartTimeSuggestions}. Phương thức này sinh ra mười hai khung giờ liên tiếp bắt đầu từ thời điểm hiện tại, mỗi khung cách nhau một giờ, rồi với mỗi khung giờ thực hiện ba bước đánh giá.

Bước thứ nhất là tra cứu lịch sử đặt chỗ của trạm trong khung giờ đó thông qua truy vấn SQL trực tiếp đến bảng chứa các bản ghi đặt chỗ, lấy số lượng đặt chỗ đã hoàn thành trong quá khứ. Hệ thống ước tính mức tải dự kiến cho khung giờ đó, rồi quy đổi thành điểm phạt theo công thức nhân số đặt chỗ với sáu, lấy min với 30 để tránh phạt quá nặng cho các giờ rất cao điểm. Bước thứ hai là ước tính thời gian di chuyển tới trạm dựa trên khoảng cách mà người dùng cung cấp chia cho tốc độ trung bình (mặc định 30 km/h khi không có dữ liệu), rồi quy đổi thành điểm phạt nhẹ bằng 0,2 lần số phút di chuyển. Bước thứ ba là tính điểm thưởng dựa trên sự gần với khung giờ ưa thích của người dùng. Khung giờ ưa thích được xác định bằng cách truy vấn lịch sử đặt chỗ của chính người dùng đó, lấy giờ trung bình trong ngày mà họ hay sạc, mặc định là 19 giờ nếu người dùng chưa có lịch sử. Điểm thưởng được tính bằng max của 0 và 8 trừ khoảng cách thời gian theo giờ giữa khung giờ đang xét và khung giờ ưa thích, với khoảng cách được tính theo vòng tròn 24 giờ để xử lý đúng trường hợp khung giờ ưa thích gần cuối ngày.

Điểm tổng hợp cho mỗi khung giờ bằng tối đa của 0 và 100 trừ điểm phạt tải, trừ điểm phạt di chuyển, trừ điểm phạt sạc cộng điểm thưởng ưa thích. Phía máy khách hiển thị ba khung giờ có điểm cao nhất dưới dạng danh sách, kèm theo thời điểm bắt đầu, điểm số và mức tải dự kiến, giúp người dùng chọn khung giờ phù hợp nhất với lịch trình cá nhân.

\subsection{Kết quả đạt được}

Tính năng gợi ý thời điểm sạc đã được đưa vào màn hình chi tiết kết quả gợi ý trong ứng dụng di động, cho phép người dùng nhấn vào nút "Thời điểm gợi ý" ngay sau khi nhận được danh sách trạm khuyến nghị. Đối với những trạm quen thuộc có nhiều lịch sử đặt chỗ, hệ thống có thể đưa ra gợi ý khác biệt rõ rệt giữa giờ cao điểm và giờ thấp điểm, giúp người dùng nhận ra lợi ích của việc sạc vào khung giờ ít tải. Đối với những người dùng mới chưa có lịch sử, hệ thống vẫn hoạt động dựa trên giả định khung giờ ưa thích mặc định là 19 giờ, vừa đảm bảo tính khả dụng vừa cho phép cá nhân hóa dần theo thời gian khi dữ liệu tích lũy. Cơ chế tích hợp chặt chẽ với thông số pin và phạm vi xe từ luồng gợi ý chính cũng đảm bảo rằng thời điểm sạc được gợi ý vẫn phù hợp với khả năng di chuyển tới trạm trong từng khung giờ cụ thể.

\section{Tổng hợp nhóm tính năng dẫn đường và gợi ý trạm}
\label{section:5.5}

Bốn phần trên đã trình bày chi tiết các đóng góp kỹ thuật cho nhóm tính năng dẫn đường và gợi ý trạm sạc. Đóng góp thứ nhất nằm ở việc kết hợp OSRM cho định tuyến với PostGIS cho truy vấn không gian, kèm theo cơ chế fallback bán kính hành lang theo bậc giúp tăng tỉ lệ tìm được trạm phù hợp ngay cả ở vùng có mật độ trạm thưa. Đóng góp thứ hai là hàm chấm điểm đa yếu tố, cho phép so sánh công bằng giữa các trạm có đặc điểm khác nhau về khoảng cách, công suất, số cổng rảnh, thời gian sạc dự kiến và điểm tin cậy, đồng thời phân biệt được trạm sạc, trạm đổi pin và trạm hỗn hợp thông qua các hằng số trọng số riêng. Đóng góp thứ ba là bộ lọc an toàn về mức pin, đảm bảo hệ thống không bao giờ đề xuất một trạm mà người dùng không thể tới được vì lý do năng lượng, một vấn đề đặc thù của xe điện mà các hệ thống dẫn đường truyền thống không giải quyết. Đóng góp thứ bốn là cơ chế gợi ý thời điểm sạc, kết hợp tải lịch sử của trạm với thói quen cá nhân của người dùng để vừa tối ưu thời gian chờ vừa tôn trọng thói quen sử dụng.

\begin{figure}[H]
    \centering
    \includegraphics[width=0.9\linewidth]{Figure/Figure_Placeholder.png}
    \caption{Giao diện dẫn đường và gợi ý trạm sạc trong ứng dụng di động}
    \label{fig:routing_ui}
\end{figure}

\begin{figure}[H]
    \centering
    \includegraphics[width=0.9\linewidth]{Figure/Figure_Placeholder.png}
    \caption{Màn hình gợi ý trạm sạc tối ưu theo thời gian di chuyển và sạc}
    \label{fig:recommendation_ui}
\end{figure}

Tất cả các đóng góp trên đều được cài đặt theo nguyên tắc "máy chủ làm hết, máy khách chỉ hiển thị", nhờ đó nhóm phát triển có thể điều chỉnh chiến lược gợi ý, thay đổi trọng số, hoặc thậm chí thay thế toàn bộ thuật toán mà không cần phát hành lại ứng dụng cho người dùng. Mã nguồn của các thành phần này được tổ chức rõ ràng trong ba lớp: lớp điều khiển \texttt{RoutingController} và \texttt{EvUserAiController} tiếp nhận yêu cầu từ máy khách, lớp dịch vụ \texttt{RoutingService} và \texttt{EvUserAiService} thực hiện tính toán, và lớp truy vấn \texttt{StationQueryRepositoryImpl} cùng \texttt{RecommendationQueryService} thực hiện các truy vấn không gian và chấm điểm. Kiến trúc phân lớp này giúp việc kiểm thử và bảo trì từng thành phần trở nên độc lập, đồng thời tạo điều kiện cho việc mở rộng hệ thống trong tương lai, ví dụ như bổ sung các yếu tố thời tiết, giờ vàng giờ bạc của lưới điện, hay tích hợp thêm các dịch vụ định tuyến khác bên cạnh OSRM.

\section*{Tổng kết chương}
\addcontentsline{toc}{section}{Tổng kết chương}

Chương này đã trình bày bốn đóng góp kỹ thuật chính của nhóm tính năng dẫn đường và gợi ý trạm sạc: pipeline dẫn đường nhận biết xe điện kết hợp OSRM và PostGIS, hàm chấm điểm đa yếu tố thống nhất cho cả trạm sạc và trạm đổi pin, bộ lọc an toàn về mức pin tới trạm, và cơ chế gợi ý thời điểm sạc dựa trên lịch sử. Mỗi đóng góp đều xuất phát từ một bài toán thực tế được xác định trong chương khảo sát, được giải quyết bằng giải pháp kỹ thuật phù hợp với kiến trúc đã trình bày ở chương 4, và đã được tích hợp vào ứng dụng di động. Trong chương tiếp theo, đồ án sẽ tổng kết lại những gì đã đạt được, so sánh với các sản phẩm cùng loại trên thị trường, rút ra bài học kinh nghiệm và đề xuất hướng phát triển trong tương lai.

\end{document}
